% ---------------------------------------------------------------------------
% Como se mede a qualidade de uma fronteira: NIGD+ e HVR.
%
% Os pontos e a semântica são os mesmos das figuras de indicadores da
% apresentação do MIC 2026 (lá feitas em pgfplots); aqui são redesenhados em
% TikZ puro, na paleta PuOr, para não carregar pgfplots nesta apresentação.
%
%   Y^ref (laranja) : fronteira de referência (todas as não dominadas).
%   S     (azul)    : fronteira obtida por um algoritmo.
%   y^ref (verde)   : ponto de referência de pior caso, usado para normalizar.
%
% Definições conforme a implementação em motmdsp/src/exec:
%   IGD+(A) = média, sobre z em Y^ref, de min_{a in A} ||(a - z)^+||
%   NIGD+   = IGD+(S) / IGD+({y^ref})
%   HVR     = HV(S, y^ref) / HV(Y^ref, y^ref)
%
% Define dois comandos, cada um uma tikzpicture completa:
%   \PainelNIGD   e   \PainelHVR
% ---------------------------------------------------------------------------

% Fronteira obtida S (4 pontos) e fronteira de referência Y^ref (7 pontos).
\def\FronteiraObtida{{(2,7)}, {(4,4)}, {(6,3)}, {(7,2)}}
\def\FronteiraRef{{(1,8)}, {(1.5,5.88)}, {(2.19,4.28)}, {(3.08,3.08)},
                  {(4.28,2.19)}, {(5.88,1.5)}, {(8,1)}}

% Eixos comuns aos dois painéis.
\newcommand{\PainelEixos}{%
  \draw[-Stealth, thick] (0,0) -- (10.4,0)
    node[below=2pt, font=\scriptsize, anchor=north east] {$C_{\max}$};
  \draw[-Stealth, thick] (0,0) -- (0,10.4)
    node[right=1pt, font=\scriptsize, anchor=south west] {$\sum_{j} C_{j}$};
}

% Marcadores das duas fronteiras e do ponto de pior caso.
\newcommand{\PainelPontos}{%
  \foreach \p in \FronteiraRef {\fill[mydarkorange1] \p circle (6pt);}
  \foreach \p in \FronteiraObtida {%
    \fill[mydarkblue1] \p ++(-3.6pt,-3.6pt) rectangle ++(7.2pt,7.2pt);}
  \node[star, star points=5, star point ratio=2.2, draw=mygreen, fill=mygreen,
        inner sep=2.4pt] at (9,9) {};
  \node[mygreen, font=\scriptsize, anchor=south] at (9,9.6)
    {$\mathbf{y}^{\mathrm{ref}}$};
}

% Legenda comum aos dois painéis, usada uma única vez no slide: mantê-la
% fora das figuras deixa cada painel quadrado e, portanto, maior.
\newcommand{\LegendaIndicadores}{%
  \tikz[baseline=-0.5ex]{\fill[mydarkorange1] (0,0) circle (0.42ex);}\,%
  fronteira de refer\^encia $\mathcal{Y}^{\mathrm{ref}}$
  \qquad
  \tikz[baseline=-0.5ex]{\fill[mydarkblue1] (-0.42ex,-0.42ex)
    rectangle (0.42ex,0.42ex);}\,%
  fronteira obtida $S$
  \qquad
  \tikz[baseline=-0.5ex]{\node[star, star points=5, star point ratio=2.2,
    draw=mygreen, fill=mygreen, inner sep=0.32ex] at (0,0) {};}\,%
  pior caso $\mathbf{y}^{\mathrm{ref}}$%
}

%%%%%%%%%%%%%%%%%%%%%%%%%%%%%%%%%%%%%%%%%%%%%%%%%%%%%%%%%%%%%%%%%%%%%%%%%%%%%
%% Painel 1: NIGD+
%%%%%%%%%%%%%%%%%%%%%%%%%%%%%%%%%%%%%%%%%%%%%%%%%%%%%%%%%%%%%%%%%%%%%%%%%%%%%
\newcommand{\PainelNIGD}{%
\begin{tikzpicture}[x=5.4mm, y=5.4mm, font=\sffamily]
  \PainelEixos

  % Normalização: distâncias de cada ponto de referência ao ponto de pior caso.
  \foreach \p in \FronteiraRef {\draw[mygreen, dotted, thick] \p -- (9,9);}

  % Para cada ponto da fronteira de referência, o ponto obtido mais próximo.
  \draw[mydarkblue1, dashed, very thick]
    (1,8)      -- (2,7)   (1.5,5.88) -- (2,7)
    (2.19,4.28)-- (4,4)   (3.08,3.08)-- (4,4)   (4.28,2.19) -- (4,4)
    (5.88,1.5) -- (7,2)   (8,1)      -- (7,2);

  \PainelPontos
\end{tikzpicture}}

%%%%%%%%%%%%%%%%%%%%%%%%%%%%%%%%%%%%%%%%%%%%%%%%%%%%%%%%%%%%%%%%%%%%%%%%%%%%%
%% Painel 2: HVR
%%%%%%%%%%%%%%%%%%%%%%%%%%%%%%%%%%%%%%%%%%%%%%%%%%%%%%%%%%%%%%%%%%%%%%%%%%%%%
\newcommand{\PainelHVR}{%
\begin{tikzpicture}[x=5.4mm, y=5.4mm, font=\sffamily]
  \PainelEixos

  % Região dominada pela fronteira de referência, limitada por y^ref.
  \draw[mydarkorange1, very thick, dashed, fill=mylightorange1, fill opacity=0.30]
    (1,9) -- (1,8) -- (1.5,8) -- (1.5,5.88) -- (2.19,5.88) -- (2.19,4.28)
    -- (3.08,4.28) -- (3.08,3.08) -- (4.28,3.08) -- (4.28,2.19)
    -- (5.88,2.19) -- (5.88,1.5) -- (8,1.5) -- (8,1) -- (9,1) -- (9,9) -- cycle;

  % Região dominada pela fronteira obtida, limitada por y^ref.
  \draw[mydarkblue1, very thick, fill=mylightblue1, fill opacity=0.55]
    (2,9) -- (2,7) -- (4,7) -- (4,4) -- (6,4) -- (6,3) -- (7,3) -- (7,2)
    -- (9,2) -- (9,9) -- cycle;

  \PainelPontos
\end{tikzpicture}}
